\section{Safety}
\label{sec:safety}

Safety is a critical aspect of embedded control systems, particularly in applications such as aerial robotics where system failures can lead to physical damage or hazardous situations. A robust system must be capable of detecting abnormal conditions and responding in a controlled manner to prevent escalation.

\subsection{Failure Modes}

Embedded flight systems can encounter various types of failures, including:

\begin{itemize}
    \item \textbf{Sensor Failures}:  
    Faulty or inconsistent sensor readings due to hardware malfunction, noise, or environmental interference.

    \item \textbf{Actuator Failures}:  
    Loss of control over motors or servos, which can directly affect system stability.

    \item \textbf{Communication Loss}:  
    Interruption of communication with remote controllers or ground stations.

    \item \textbf{Software Faults}:  
    Unexpected behavior due to bugs, race conditions, or incorrect state estimation.
\end{itemize}

Identifying these failure modes is the first step toward designing effective safety mechanisms.

\subsection{Fault Detection}

To ensure safe operation, the system must continuously monitor its internal state and sensor inputs. Fault detection mechanisms may include:

\begin{itemize}
    \item Threshold-based checks on sensor values and system states
    \item Consistency checks between multiple sensors (e.g., comparing accelerometer and gyroscope estimates)
    \item Monitoring of communication timeouts
    \item Detection of actuator saturation or abnormal control outputs
\end{itemize}

Early detection of anomalies allows the system to take corrective action before instability occurs.

\subsection{Failsafe Mechanisms}

Failsafe mechanisms define how the system responds when a fault is detected. These responses are designed to minimize risk and maintain system safety.

Common failsafe strategies include:

\begin{itemize}
    \item \textbf{Disarm}:  
    Immediately disable actuators to prevent unintended motion.

    \item \textbf{Hover / Stabilize}:  
    Maintain a stable state using available sensor data when partial failures occur.

    \item \textbf{Return-to-Home (RTH)}:  
    Navigate the system back to a predefined safe location in case of communication loss.

    \item \textbf{Controlled Shutdown}:  
    Gradually reduce system activity to avoid abrupt transitions.
\end{itemize}

The choice of failsafe strategy depends on system capabilities and the nature of the detected fault.

\subsection{Design Considerations}

Designing safety mechanisms involves balancing responsiveness and reliability. Key considerations include:

\begin{itemize}
    \item \textbf{Latency of Detection}:  
    Faults must be detected quickly to prevent escalation.

    \item \textbf{False Positives}:  
    Overly sensitive detection may trigger unnecessary failsafe actions.

    \item \textbf{System Redundancy}:  
    Multiple sensors or fallback mechanisms can improve reliability.

    \item \textbf{Graceful Degradation}:  
    The system should maintain partial functionality when possible, rather than failing abruptly.
\end{itemize}

A well-designed safety framework ensures that the system can handle unexpected conditions in a predictable and controlled manner, forming an essential component of reliable embedded systems.